% SPDX-License-Identifier: Apache-2.0
\section{An Array of Ports and a Loop}
\label{sec:x-loop}

A unit of a count used to be text a macro wrote once per count:
\code{router!} writes \code{Router2} to \code{Router8}, and the
largest it writes is the largest there is. A side of \code{run} may
now be an array of ports, \code{ins: [Rx<U<8>>; N]}, and the body may
go over it with \code{for i in 0..N} (issue 500).

\code{N} is a const parameter, which the macro, reading tokens, cannot
see. So the loop is unrolled when \code{lowered} runs rather than when
the macro expands: the loop becomes a Rust \code{for} around the
statements it adds, and \code{ins[i]} is the port \code{ins\_0},
\code{ins\_1} and so on as it goes round. The array is a struct of
ports whose fields are its indices, so its ports are found the way a
\code{LitePort}'s are. Inside the loop, and for a \code{let mut}, a
\code{let} is a wire of its own on every turn, and assigning to a
\code{let mut} makes the next one, which is how \code{found} and
\code{pick} carry what the loop has seen from one input to the next.
Each turn's wires are written once, so the netlist grows with \code{N}
and not with its square.

\code{Merge<N>} forwards one word a cycle from the lowest input that
has one. The run offers words on three inputs at once and asserts that
they come out lowest input first and none is lost, and the netlist of
\code{Merge<3>} is checked against the run under both simulators. The
same source lowered at five has five input ports, which the run
prints. The AXI arbiter, \code{Arbiter<N, ..>}, is written this
way too: its hosts are three arrays of ports in and two out, and its
grant is a loop that carries the first host it has found. A \code{for} under an \code{if} is refused: put the \code{if}
inside the \code{for}. The index is written \code{ins[i]} because that
is what the lowering reads, so Clippy's wish for an iterator is
allowed with the reason beside it.

\begin{figure}[htbp]
\centering
\input{loop_timing.tex}
\caption{Three inputs into one: while more than one input has a word,
the lowest goes first.}
\label{fig:x-loop}
\end{figure}

\lstinputlisting[caption={\code{ex\_loop.rs}. Run with \code{bazel run //lib/examples:ex\_loop}.},label={lst:x-loop}]{lib/examples/ex_loop.rs}

\lstinputlisting[style=ascii,caption={What \code{ex\_loop} printed when the build ran it: the words in the order they came out, the lowering of \code{Merge<3>}, and the ports of \code{Merge<5>}.},label={lst:x-loop-out}]{out_loop.txt}
